\chapter{Credit losses and the bank balance sheet}
\label{ch:banking}

The funding and default distinctions become central in a bank. A loan creates
an asset, customer payments affect settlement liquidity, and a forecast of
credit loss affects an allowance before cash is lost. The first case connects
funding to a simple allowance; the later dated loan schedule separates credit
stages, interest recognition and eventual recovery.

\section{Banks: accounting capital is not regulatory capital}

A bank package needs instrument and cohort states for loans, securities,
deposits, wholesale funding, repos, derivatives, collateral and commitments.
For interest-rate analysis, currency and repricing buckets matter alongside
remaining maturity. For credit analysis, exposure, default state, loss severity,
recoveries, guarantees and their dependencies matter. For liquidity analysis,
deposit behaviour, collateral eligibility, encumbrance and payment timing matter.
These states can be represented at an explicitly chosen aggregation level;
the public record rarely identifies every contract.

\subsection{A reported group boundary}

JPMorgan Chase's filings make the boundary between a bank, its parent and its
securities issuers concrete. In the 2022--2025 reports, Item~1 identifies
JPMorgan Chase Bank, N.A. as the principal bank subsidiary and describes an
intermediate holding company, JPMorgan Chase Holdings LLC (the IHC). The
reports also describe J.P.~Morgan Securities LLC as a principal non-bank
subsidiary and place J.P.~Morgan Securities plc and J.P.~Morgan SE beneath the
bank. These are reported relationships for the stated dates, not a complete
legal-entity graph \citep{jpmorgan2022tenk,jpmorgan2023tenk,jpmorgan2024tenk,jpmorgan2025tenk}.

The attached Exhibit~21 tables qualify their scope: they list significant
legal-entity subsidiaries, some parents of significant subsidiaries and certain
entities designated for resolution planning, while excluding subsidiaries that
would not be significant in aggregate under the stated SEC definition. Their
indentation shows how the issuer presents the rows, but it does not supply
ownership percentages, tax residence, functional currency or unrestricted
transfer rights. A model therefore keeps a reported display relationship
separate from a confirmed control edge.

Exhibit~22.2 adds a different edge. It names JPMorgan Chase Financial Company
LLC as issuer of Global Medium-Term Notes, Series~A, and JPMorgan Chase \& Co.
as guarantor. Financial Company LLC does not appear in those qualified
Exhibit~21 lists. The 2022--2023 text describes Series~A notes as currently
offered and sold under one cited registration; the 2024--2025 text also says
that the issuer may issue debt securities and warrants under a later cited
registration. The descriptions say that the relevant securities are fully and
unconditionally guaranteed, but they give no outstanding principal, guarantee
valuation or accounting recognition. They also do not establish that a warrant
was issued or that an earlier note was redeemed.

The same reports state that the parent depends on subsidiary cash flows to pay
its securities. The bank's dividends are subject to regulatory restrictions,
and the IHC cannot pay dividends or extend credit to the parent when stated
capital or liquidity thresholds, or management and board limits, apply. This
means that consolidated cash, a subsidiary's assets and the parent's
distributable funding are different state variables. The example gives the
simulator a place to record the legal entity, guarantee and conditional funding
path before it assigns a cash movement; it does not infer a balance or a
probability of default.

Note~33 makes the funding distinction measurable for 2025. Its parent-only
income statement reports net income of US\$57{,}048m. The parent-only cash-flow
statement subtracts US\$71{,}163m of subsidiary and affiliate net income,
reports a US\$14{,}115m parent company net loss in that reconciliation, then
adds US\$50{,}000m of cash dividends from subsidiaries and US\$8{,}583m of
other operating adjustments to reach US\$44{,}468m of operating cash. Neither
the income figure nor the dividend alone is the period's ending cash.

At 31~December 2025, the parent reports US\$39m of cash and due from banks and
US\$9{,}751m of deposits with banking subsidiaries, together US\$9{,}790m.
The cash-flow footnote says this measure includes US\$8.0bn of intercompany
time deposits. The same parent-only balance sheet reports US\$669{,}449m of
equity-method investments in bank and bank-holding subsidiaries; that asset is
not parent cash. Operating cash of US\$44{,}468m and investing cash of
US\$34m, less US\$44{,}512m used in financing, leave closing cash and
deposits US\$10m below their opening balance.
These dated amounts distinguish accounting income, existing cash and
conditional subsidiary transfers. A future debt-service calculation still
needs payment dates, permissible transfers, deposit terms and funding actions
\citep[pp.~19, 310--311]{jpmorgan2025tenk}.

Note~26 supplies a dated transfer capacity, which differs from both cash held
and cash already received. It reports that, at 1~January 2026, JPMorgan Chase
Bank, N.A. could pay approximately US\$30.1bn in dividends to the parent
without prior approval from its relevant banking regulators. That is a
capacity of the named bank at that date; the US\$50{,}000m above is cash
received from subsidiaries during 2025. The note also describes regulators'
power to prohibit unsafe or unsound dividends and the IHC's separate capital,
liquidity and management restrictions. It says the bank's 2026 earnings will
supplement dividend capacity, without estimating those earnings. Turning the
reported capacity into a future cash receipt still requires a dividend decision,
payment date and the applicable conditions \citep[pp.~4, 292]{jpmorgan2025tenk}.

The same note reports US\$29.0bn of consolidated restricted cash at
31~December 2025: US\$19.4bn segregated for securities and cleared-derivative
customers, plus US\$9.6bn of non-US central-bank reserves and cash held for
other general purposes. Its balance-sheet footnote allocates the total to
US\$27.8bn of deposits with banks and US\$1.2bn of cash and due from banks.
Those classifications describe the same restricted-cash total. Subtracting it
from the parent's US\$9{,}790m would mix consolidated and parent-only
perimeters and cannot determine the parent's unrestricted cash. Other pledged
cash and securities are disclosed separately; their category overlap must be
established before adding them to a restricted-resource total
\citep[p.~292]{jpmorgan2025tenk}.


For a future cash dividend, the model needs the bank's declaration as well as
the eventual payment. The OCC's national-bank dividend rules, retrieved for
1~January 2026, use the declaration date to determine compliance. They define
retained net income as net income less all dividends declared in the period.
Absent OCC approval under the earnings provision, current-year declarations, including
the proposal, are compared with current-year net income to date plus the
previous two years' retained net income, less specified required transfers.
The current year is the calendar year of declaration. A cash-flow statement
of dividends paid therefore cannot supply the declaration history by itself
\citep[12~CFR~5.61--5.64]{occnationaldividends2026}.

The preceding two years are not always enough. Section~5.64(c)(2) provides
conditional offsets for retained losses caused by dividends exceeding a
single year's net income. Those offsets can require the third and fourth
preceding years and follow a specified order; other earnings deficits do not
qualify for that adjustment. An implementer must preserve the cause of the
loss and the allocation of offsets before computing the earnings restriction.
Capital, undivided-profit and supervisory conditions remain separate: an
earnings calculation alone cannot establish a permissible transfer
\citep[12~CFR~5.63--5.65]{occnationaldividends2026}.

Selecting the applicable capital rule also requires more than a calendar
lookup. The final leverage rule published on 1~December 2025 is effective
1~April 2026 but permits early adoption from 1~January
\citep[p.~55248]{eslrfinal2025}. JPMorgan's 2025
Form~10-K reports that the Firm elected early adoption effective 1~January
2026 \citep[p.~91]{jpmorgan2025tenk}.
The rule's publication, general effective date, permitted election date and
the issuer's reported election answer different questions. The later filing
provides evidence about the election; it does not establish that the election
was publicly known on 1~January. A historical information query must retain
that distinction even when reconstructing an earlier regulatory state.

The amendment removes the additional 6\% supplementary-leverage condition
from the affected national banks' well-capitalized classification and adds
separate leverage-buffer payout provisions. That classification change does
not remove the distinct undercapitalization test or establish unrestricted
distributions \citep[pp.~55287--55288]{eslrfinal2025}. The reported Firm election
also does not supply every bank-specific application fact. A permissible
dividend calculation still needs the relevant capital definitions, distribution
limits, supervisory conditions and bank-level inputs. The US\$30.1bn remains
the issuer's dated capacity disclosure, rather than a recomputed legal limit.


One of those inputs, eligible retained income, differs from the
retained net income used in the national-bank earnings rule. Part~3 takes the
greater of two amounts: Call Report net income for the four preceding calendar
quarters, net of distributions and associated tax effects not already reflected
in that income; and the average quarterly net income over those four quarters.
Its maximum payout amount for the current quarter multiplies this eligible
income by the applicable payout ratio. The relevant capital buffers use the
previous quarter's closing ratios, so the declaration, income window and
capital-measurement date must remain distinct
\citep[12~CFR~3.11(a)(2)--(4)]{occcapitalpayout2026}.

The transactions being limited also extend beyond a cash dividend. Part~3's
distribution definition includes dividend declarations or payments on Tier~1
capital instruments and specified capital repurchases or redemptions, with
conditional exceptions for qualifying replacements
\citep[12~CFR~3.2]{occcapitaldefinitions2026}. The quarterly limit also
covers discretionary bonus payments and the creation of obligations to make
the specified payments. A bank's event history must therefore identify the
instrument, declaration, obligation and settlement, then reconcile their
regulatory treatment; a cash-dividend total cannot substitute for that history
\citep[12~CFR~3.11(a)(4)]{occcapitalpayout2026}.

A hypothetical bank in the first quarter of a year illustrates the income
distinction. Suppose each of the preceding four quarters has Call Report net
income of US\$10m. Assume its only distributions in that preceding calendar
year were US\$38m of declared dividends, with no associated tax adjustment.
Part~5 retained net income for that year is then US\$2m: US\$40m less
US\$38m. Part~3 eligible retained income is US\$10m, the greater of that
US\$2m remainder and the US\$10m quarterly average. These are two calculations
for the same hypothetical history, not interchangeable names for one balance.

Now suppose the bank falls within the amended rule's covered population and
applies that rule. Take its previous-quarter capital conservation buffer as
1.5\%, its applicable countercyclical amount as zero, its leverage buffer
standard as 1\%, and its measured leverage buffer as 0.2\%. These are supplied
example inputs, rather than estimates from a public balance sheet. The rule's Table~1
gives a 40\% payout ratio, which alone would produce US\$4m from the
US\$10m eligible income. Table~2 gives zero because the leverage buffer is
at most one quarter of its standard. The required lower ratio therefore gives
a zero table-based maximum, absent the specified OCC permission. Conversely,
a table entry saying that no payout-ratio limitation applies does not mean a
100\% ratio or permission under every other rule
\citep[12~CFR~3.11(a)(4), (c), Tables~1--2]{occcapitalpayout2026}.


Replacing the example's supplied leverage-buffer standard with an observed
requirement needs a specific surcharge. For a bank within the amended rule's
covered population, the standard is the lesser of 1\% and, if applicable,
half of the controlling global systemically important holding company's most
recent required Method~1 surcharge, subject to the specified effective-date
provisions. The holding company's overall surcharge instead takes the greater
of its Method~1 and Method~2 surcharges. The overall figure alone does not
identify the input this bank rule requires.
Annual determination and the date an adjusted surcharge takes effect also
remain separate
\citep[12~CFR~3.11(a)(1)(v)]{occcapitalpayout2026}
\citep[12~CFR~217.403(a)--(d)]{fedgsibsurcharge2026}.

Method~1 compares twelve reported systemic indicators with published global
amounts, applies their specified weights and caps the substitutability
category's contribution at 100 basis points
\citep[12~CFR~217.404]{fedgsibscore2026}.
The denominators have their own measurement and publication dates. The
Federal Reserve's notice published on 16~December~2025 supplies the 2025
denominators using year-end 2024 global data, converted from euros into
dollars at the stated 31~December~2024 rate. The preceding notice, published
on 18~December~2024, uses year-end 2023 data
\citep[pp.~58245--58246]{fedgsibdenominators2025}
\citep[pp.~102908--102909]{fedgsibdenominators2024}.
The release-year label therefore does not identify the underlying observation
year or the date a bank's capital requirement changes. The published amounts
are in dollars; reported indicator values must be brought to the same unit
before division.

Reporting precision adds another constraint. The retained FR~Y-15 general
instructions require dollar amounts in thousands and permit rounding to the
nearest million, with adjusted detail so the components add to the totals.
Its payment-activity instructions instead say payment totals are not rounded,
specify accuracy according to their magnitude and allow known overestimates
when precise totals are unavailable. A reconstruction cannot treat every
reported input as an exact value or give all components independent rounding
intervals. These dollar-reporting rules also leave the convention for a
fractional basis-point score unresolved
\citep[GEN-3--4 and C-1; physical pp.~7--8, 25]{fry15instructionsretained2026}.

\begin{samepage}
The report's quarter-end date differs from its submission deadline and public
availability; amendments, later resubmissions and confidentiality can change
what a historical reader could know. The retained instruction file itself
combines a September~2021 cover with a January~2026 reporting-scope page.
Applying it consequently requires the relevant page and reporting vintage,
alongside the applicable rule and institution facts. These sources explain
the input requirements; they do not yet supply a verified bank-specific score
or a complete reporting-rule history
\citep[cover and GEN-1--5; physical pp.~1, 5--9]{fry15instructionsretained2026}.
\par
\end{samepage}


Debt amounts add a second distinction. Note~20 puts US\$253{,}039m of
long-term-debt carrying value under Parent company for 2025, while Note~33's
parent-only balance sheet shows US\$254{,}633m. The US\$1{,}594m difference
equals Note~20's separately displayed junior-subordinated-debt subtotal. The
same comparison gives US\$1{,}603m in 2024. This exact two-year arithmetic
match does not establish how that row relates to the parent-only balance,
or which legal entity owes each instrument. Its accounting and
contractual relation still needs source support.

Timing requires a different measure. Note~20 places US\$14{,}298m of parent
carrying value in its under-one-year bucket; Note~33 reports US\$14.3bn of
parent long-term debt contractually maturing in 2026. They align at the
footnote's displayed precision, but carrying value includes premiums,
discounts, costs and valuation adjustments rather than simply principal
payable. For the separately reported consolidated zero-coupon population,
Note~20 gives US\$18.1bn of carrying value against US\$57.6bn of aggregate
principal at the notes' respective maturities. The firm's US\$435{,}206m
consolidated long-term-debt total and US\$47.6bn of subsidiary debt and
structured notes guaranteed by the parent are also separate from the parent's
own debt balance; the latter guarantee amount supplies neither a liability
valuation nor a payment forecast. A dated parent debt-service path still needs
instrument terms, cash dates and permissible funding transfers
\citep[pp.~281--282, 310--311]{jpmorgan2025tenk}.

\newpage

The complete Form~10-Ks describe Series~A as the issuer's direct, unsecured
general obligations, ranking equally with its other unsecured and
unsubordinated debt. Claims under the parent guarantee generally rank behind
those of subsidiary creditors. Their summary of the 2016~Indenture requires
the parent, on demand after issuer nonpayment, to make full and punctual
payment of principal, interest and other amounts due: a guarantee of payment,
not collection
(Sections~14.01--14.02). It excludes amounts the issuer is excused from paying
or beyond those actually due. Issuer defenses and counterclaims carry over,
except those arising solely from insolvency or bankruptcy status. These
descriptions refer to the Indenture and applicable prospectus supplements;
they specify reported terms but do not supply a current balance, valuation or
accounting recognition.

The complete filings also describe one particular Series~A population: the
Callable Fixed Rate Notes due 10 June 2032. They report US\$5,000,000 of
aggregate principal outstanding at 31 December 2022 in the 2022 filing, at
31 December 2023 in the 2023 filing, and at 31 December 2024 in both the 2024
and 2025 filings. The last report repeats a 2024 as-of date; this passage
supplies no 2025 year-end balance. The notes were issued on 10 June 2021 and bear 2.60\%
annual interest payable semiannually on 10 June and 10 December. They mature
on 10 June 2032 unless an earlier contractual outcome applies. The issuer
may redeem all, but not part, on 10 June 2031 at principal plus accrued and
unpaid interest. Upon acceleration after a continuing event of default, the
stated amount is US\$1,000 per US\$1,000 principal plus accrued and unpaid
interest. The descriptions report NYSE listing under JPM/32, DTC global
book-entry issuance, and the possibility that a reopening increases aggregate
principal outstanding \citep{jpmorgan2022tenk,jpmorgan2023tenk,jpmorgan2024tenk,jpmorgan2025tenk}.

These dated amounts are reported principal, not carrying or fair value. Nor do
they establish the accounting classification or value of the parent guarantee,
the entire Series~A population, or an unchanged balance between the stated
dates. Each filing's Exhibit~4.4 identifies the 19 February 2016 Indenture
and incorporates it by reference to Exhibit~4(a)(7) of Form S-3 file
333-209682, filed 24 February 2016. The complete filings summarize that
instrument; its text and applicable supplements or amendments require their
own inspection before being used as complete legal terms.

\newpage
The filings then describe general Series~A terms that may apply to securities
including the Callable Notes. In their summary, unpaid interest must continue
for 30 days to become an event of default; principal and other amounts default
when due. An issuer covenant breach has a 90-day period after the specified
written notice. For listed continuing payment, issuer-covenant or guarantee
events, the trustee or holders of at least 25\% of the aggregate principal
then outstanding for the relevant series may, by written notice to the issuer
and guarantor, declare principal and accrued interest due. Specified continuing
issuer insolvency events accelerate the debt automatically without a declaration. Guarantor insolvency
alone is excluded and does not authorize acceleration under this description,
while a payment default by the guarantor remains a separate possibility
\citep{jpmorgan2022tenk,jpmorgan2023tenk,jpmorgan2024tenk,jpmorgan2025tenk}.
A history must therefore keep the failure, notice, elapsed time, actor,
declaration and changed due schedule distinct, rather than record one
group-wide default flag.

Collective remedy control also differs from an individual's right to payment.
The summaries condition an ordinary holder remedy action on prior trustee
notice, a request by holders meeting the stated 25\% threshold, 60 days of
trustee inaction, indemnity and no inconsistent majority direction. They
separately protect each holder's right to receive and sue for payment when due
against impairment without that holder's consent. The summary permits certain
amendments without holder consent under Section~8.01. For consent-based
modifications under Section~8.02, it specifies a majority in aggregate
principal amount of each affected outstanding series, whereas
listed changes to maturity, amounts, currency, payment method, redemption,
payment suit or specified holder repurchase rights, voting threshold or adverse guarantee terms require each
affected holder's consent. These thresholds attach to different actions and
voting populations, not a generic creditor vote. The 2022 default-and-waiver
subsection omits the trustee-conflict paragraph that the 2023--2025 versions
place there. The 2022 filing discusses a potential default, trustee conflict
and possible resignation later under its trustee heading. This change in
placement does not establish a contract amendment. The original Indenture and
series-specific supplements remain uninspected
\citep{jpmorgan2022tenk,jpmorgan2023tenk,jpmorgan2024tenk,jpmorgan2025tenk}.

\newpage
Reported succession terms give another reason to keep the obligor separate
from the debt claim. Under the general Series~A summary, a successor to the
issuer in a covered merger, consolidation or transfer of all or substantially
all assets must be a U.S.\ corporation or limited liability company that
expressly assumes due payment and the issuer's covenants. A successor to the
guarantor must instead be a U.S.\ corporation that assumes the full guarantee
and the guarantor's covenants. The continuing issuer or guarantor is the other
stated route. Each party's route also requires that no default or event that
would become a default with notice or time has occurred or would be continuing
immediately after the transaction. The descriptions deem certain
resolution-related transfers of the guarantor's material assets to meet the
all-or-substantially-all category; they do not define a numeric percentage
\citep{jpmorgan2022tenk,jpmorgan2023tenk,jpmorgan2024tenk,jpmorgan2025tenk}.

The same summaries say a recapitalization, change of control or highly
leveraged transaction alone receives no additional holder protection from
these merger provisions; a particular supplement may provide a put or higher
interest. The model therefore needs a dated successor, assumed payment and
covenant duties, separate issuer/guarantor conditions and any evidenced
series-specific protection. Succession itself does not prove repayment,
accounting derecognition or unrestricted access to the group's cash. These
are reported terms that may apply, pending inspection of the original
Indenture and applicable supplements.

Loan origination can debit loans and credit a customer deposit in the bank's
book. Subsequent transfer to another bank produces settlement flows. Treating
all loan growth as an immediate cash outflow would therefore encode a specific
funding assumption, not a universal banking identity. A system-wide model must
also link counterparties and settlement assets; single-company simulation
records external counterparties without claiming to close the whole economy.

\newpage
An allowance calculation is a dated measurement based on forecast information.
For a single horizon, let EAD be a fixed exposure at default, PD the probability
of default, and LGD the mean loss fraction conditional on default. Ignoring
discounting, the illustrative expected loss in currency units is
\[
 L=EAD\times PD\times LGD.
\]
With exposure 100, default probability 0.02 and loss given default 0.4, it
produces 0.8. Changing those inputs to 0.04 and 0.5 produces 2, so the allowance
increase is 1.2 if there are no other movements. This calculation is a
demonstration model. A standard-specific expected-credit-loss implementation
must determine horizon, staging or other classification, discounting,
forward-looking information, write-offs and purchased-credit-impaired or
other special cases. The simple formula is not an IFRS 9 or US GAAP compliance
implementation.

The exposure entering an allowance calculation also needs a defined
population. JPMorgan's 2025 report gives four different amounts for
lending-related commitments at 31 December 2025:
\begin{center}
\begin{tabularx}{\linewidth}{@{}Xr@{}}
\toprule
Reported quantity & US\$ millions \\
\midrule
Contractual lending-related amount & 1,817,307 \\
Commitments in the allowance-methodology table & 603,947 \\
Allowance for lending-related commitments & 5,071 \\
Lending-related carrying value & 6,278 \\
\bottomrule
\end{tabularx}
\end{center}
Notes 13 and 28 define these populations and measurements separately
\citep[pp.~258--262, 295--297]{jpmorgan2025tenk}. The contractual total
contains consumer and wholesale commitments; the allowance table excludes
specified populations. The commitment allowance is reported in accounts
payable and other liabilities. Carrying value can also contain guarantee
liabilities, fair-value measurements and deferred fees or purchase adjustments,
as applicable. For standby letters of credit and other financial guarantees,
the report explicitly reconciles a US\$175m allowance and a US\$441m guarantee
liability to US\$616m of carrying value. Treating either the contractual total
or carrying value as the allowance would compute the wrong quantity.

The report also distinguishes contractual capacity from expected funding.
Customers may leave facilities unused, refinance them or let them expire;
some commitments can be reduced or cancelled under their contracts or law.
The reported policy records no allowance for future draws on unconditionally
cancellable commitments. That exclusion concerns the loss measurement;
a liquidity scenario still needs the conditions under which a customer can
draw, the bank can cancel, and a payment requires settlement funds. Those
conditions and their dependence on borrower distress cannot be recovered
from a commitment total. Comparisons also need the same reporting basis:
specified wholesale amounts are net of risk participations here but gross
in Federal Reserve filings. Parent guarantees of subsidiary debt are
expressly outside the Note~28 table
\citep[pp.~258, 295--296, 300]{jpmorgan2025tenk}.

The historical FASB amendment ASU~2016-13 helps explain these distinctions.
Its paragraph 326-20-30-11 measures off-balance expected credit losses over
the contractual period of a present obligation to extend credit, except for
an obligation the issuer may cancel unconditionally. For a covered commitment,
the estimate considers both whether funding will occur and the expected loss
on commitments that do fund. Example~10 keeps funded credit-card loans within
loss measurement while excluding an allowance on unconditionally cancellable
undrawn credit, even though the bank has sometimes advanced funds before
detecting default. Paragraph 460-10-30-5 separately requires a contingent
expected-loss liability and a noncontingent fair-value guarantee liability
for covered financial guarantees. These are historical amendment provisions,
not verified current Codification; applying them to the 2025 report requires
the actual contracts and later rule history
\citep[PDF pp.~2, 108, 118, 144, 173--176]{fasb2016asu}.

A commitment to acquire loans makes the time distinction especially clear.
The report dates the Apple Card forward purchase commitment to 30 December
2025, its announcement to 7 January 2026, and expected closing to
approximately 24 months later. The commitment disclosure includes about
US\$104bn of estimated credit exposure at closing, including about US\$23bn
of estimated drawn loans, while the reported transaction allowance is
US\$2.2bn \citep[pp.~47, 260, 296]{jpmorgan2025tenk}. The expected drawn
loans at that future closing are not loans already owned at the 2025 reporting
date. A reconstruction therefore retains the forward commitment, the future
portfolio estimate and the dated allowance as different facts. A forecast
must supply closing and subsequent draw assumptions before generating loan
and settlement events. These reported policies and amounts support the
distinctions; their application to a new calculation still requires the
operative accounting rules and actual contract facts.

Regulatory capital is computed in a separate prudential layer with its own
perimeter, eligibility, deductions, risk weights and transition provisions.
Accounting equity feeds a reconciliation to capital; it is not automatically
Common Equity Tier 1. Liquidity and capital constraints influence the
management policy, which in turn generates asset sales, funding or capital
events. A breached ratio can be an economically meaningful scenario outcome.
No hidden capital injection is permitted to repair it.

The United States national-bank rules show why that reconciliation needs
records for each instrument. Under the common-stock criteria in the retained
January and April 2026 texts, classification as equity under GAAP is only one
condition. The instrument must also satisfy conditions concerning payment and
direct issuance by the bank, priority in insolvency, loss absorption, permanence,
dividend discretion, purchase financing and guarantees. The rules distinguish
common equity Tier~1, additional Tier~1 and Tier~2, each with its own eligible
elements and deductions. A security's name and its balance-sheet category
therefore cannot supply its regulatory treatment
\citep[12~CFR~3.20(a)--(d)]{occcapitaleligibility2026}.

Regulatory recognition can also change while a debt remains outstanding.
The general Tier~2 instrument criteria require an original maturity of at
least five years. At the beginning of each of its last five years, the amount
eligible for Tier~2 falls by 20\% of the original amount, net of redemptions;
the instrument is excluded when less than one year remains. Calls and early
redemptions have further conditions, including OCC approval. This reduction
in eligible capital does not itself reduce the accounting debt. The contract,
issuance and redemption history must support both amounts separately
\citep[12~CFR~3.20(d)(1), including footnote~7]{occcapitaleligibility2026}.

Deductions require similarly specific records. Section~3.22 deducts goodwill
and specified other assets from common equity Tier~1, with separate rules for
their tax effects. Its general rule permits an associated deferred tax
liability to reduce an asset deduction only if that liability would be
extinguished when the asset is impaired or derecognized. The same liability
cannot offset two assets; separate conditions govern offsets against deferred
tax assets, and the chosen netting treatment must be consistent across
reporting periods unless the OCC approves a change. Accumulated other
comprehensive income has its own adjustments and election conditions,
including limits on which institutions may opt out and when they must elect.
Reconstructing capital consequently needs asset-specific tax links, instrument
facts, institution classification and the actual election history. The reported
equity and aggregate deferred-tax balances do not determine that reconstruction
\citep[12~CFR~3.22(a), (b), (e)]{occcapitaldeductions2026}.


\section{A bank's cash, allowance and capital in one path}

The bank reference state contains settlement cash 30, gross loans 100,
allowance 0.8, deposits 110 and equity 19.2. Originating a loan of 20 first
debits loans and credits the borrower's deposit. A subsequent transfer of 15
to another bank debits deposits and credits settlement cash. Loan origination
and external cash use are consequently distinct events. Receive loan interest
6 in cash and credit deposit interest 2 to customer balances. Remeasure the
allowance to 2.4, posting an additional loss of 1.6. The resulting state is
cash 21, gross loans 120, allowance 2.4, deposits 117 and equity 21.6.
Net assets are $21+120-2.4=138.6=117+21.6$.

Now 25 of deposits runs off. The declared funding policy uses cash first,
then a committed facility with limit 10. It draws 4 and settles 25, leaving
cash zero, deposits 92 and wholesale debt 4. A runoff of 40 instead leaves
9 unpaid even after the maximum facility draw. The simulator records default
and retains the unpaid deposit obligation. It does not assume unlimited
central-bank or shareholder funding.

The prudential calculation is a separate, illustrative policy. Deduct 2 from
accounting equity to obtain capital 19.6; apply a 75\% risk weight to gross
loans for risk-weighted assets 90. The resulting ratio is $19.6/90$.
Neither the deduction nor the weight is a Basel default. Actual prudential
packages require jurisdiction, exposure class, eligibility, buffers and dated
rules, with a reconciliation from the accounting book. The structure follows
the need to model capital and liquidity jointly in
\citet[Section 2]{behn2019}, without transferring that paper's estimated
preferences or shock parameters to this bank.

For a richer credit adapter, partition borrower cohorts by currency, vintage,
credit state, remaining maturity, collateral and measurement category. At date
$t$, supply scenario probabilities, conditional default-time probabilities,
exposure paths, conditional loss severity and discount factors. Its expected
loss calculation has the declared form
\[
 L_t=\sum_s p_s\sum_{k\in H}
  d_{t,k}\Pr(\tau=k\mid s,\mathcal F_t)
  \E[\mathrm{EAD}_k\mathrm{LGD}_k\mid\tau=k,s,\mathcal F_t].
\]
The conditional product expectation retains dependence between exposure and
severity. The horizon $H$ concerns eligible default times; associated cash
shortfalls may continue later. The rule package decides twelve-month versus
lifetime default scope, credit-state transfers, interest basis and write-off
recognition. Changes in allowance must be attributed to origination,
derecognition, transfers, revised expectations, currency and write-offs.
Contract and prudential support can be specified now, while full IFRS~9/CECL/
ASBE classification and paragraph-level validation remain later obligations.

\hypertarget{state-and-event-order}{%
\section{State and event order}\label{method-bank-methods-state-and-event-order}}

The dated credit methods use the selected amortised-cost, interest and impairment provisions of AASB 9 \citep{aasb9dated}. The following loan has different inputs from the earlier liquidity example, so its allowance and final equity are calculated separately.

BankState contains cash/reserves by settlement venue and currency; deposits by customer, access rights and maturity; assets by contract, gross carrying amount, allowance, original EIR, accrued cash flows, collateral and credit state; undrawn commitments; funding lines and covenants; accounting equity; regulatory mappings; and evidence/information views. LoanState separately retains origination credit risk and current risk, SICR finding, credit-impaired/POCI status, forbearance, modification, write-off and recovery history. Aggregating by today's loan balance alone loses both staging evidence and later response to shocks.

Each time step first admits information known then, forms permitted credit and liquidity forecasts, chooses behavior under that information, and schedules due events. Accrual, actual settlement, impairment measurement and regulatory queries have separate business identities. A deposit received from outside debits cash/reserves and credits deposit liability. A loan credited to the borrower's deposit debits loans and credits deposits without initially transferring reserves; a transfer to another bank then debits deposits and credits reserves. A loan paid directly outside instead debits loan/credits reserves. These two funding paths cannot be treated as identical cash events. Deposit interest debits interest expense/credits deposit or accrued payable; withdrawal requires settlement funds.

\hypertarget{expected-credit-loss}{%
\section{Expected credit loss}\label{method-bank-methods-expected-credit-loss}}

For each admitted scenario j and date t, calculate contractual due flows minus the cash expected to be received, including supported collateral/enforcement cash and timing; discount the shortfall under the applicable original EIR or specified exception. With mutually exclusive exhaustive scenario weights w\_j, ECL is the sum of w\_j times those present-value shortfalls. Default-time hazards must be converted to disjoint default-time probabilities; summing cumulative PDs over dates counts the same default repeatedly. Collateral recoveries cannot be credited both in LGD and as a second recovery cash flow. A shorthand PD\ensuremath{\times}LGD\ensuremath{\times}EAD is admissible only if its definitions and discounting reproduce the specified cash-shortfall calculation; it is not a universal replacement.

Expected credit loss (ECL) measures the discounted expected cash shortfall. A significant increase in credit risk (SICR) changes the general method's default horizon. Purchased or originated credit-impaired (POCI) assets follow a separate method. For the general approach, twelve-month ECL considers lifetime shortfalls arising from possible defaults during the next twelve months; it is not just unpaid cash in those twelve months. SICR uses the change in lifetime default risk since initial recognition, reasonable forward-looking information and the applicable rebuttable presumptions. Do not install a two-times-PD or 30-day rule as an unconditional universal trigger. Lifetime ECL follows a supported SICR finding; cure needs a supported reversal of that finding. Subsequently credit-impaired assets use original EIR on amortised cost for interest in subsequent periods; POCI has a separate credit-adjusted-rate and cumulative-change method.

The 2017 ASBE~22 text makes the two horizons explicit. Article~47 defines
credit loss as the present value of all contractual cash flows less expected
receipts; a late receipt can therefore create a loss even when its nominal
amount is paid in full. Article~48 limits twelve-month ECL by the date of the
possible default event, while retaining the resulting cash shortfalls over the
asset's life. Article~52 assesses changes in default risk, rather than changes
in the allowance, and Article~58 requires probability weighting, time value
and reasonable, supportable information \citep[Articles 47--48, 52, 58]{mofasbe22}.
These provisions support the following dated comparison. Article~84 states
1 January 2018, but applying the standard to a particular company also requires
its issuance schedule, subsequent amendments and adoption facts.

The general measurement interface requires scenario\_cashflows, law\_ref, information\_view, horizon\_admission, credit\_assessment and discount\_rule. An unweighted history set can produce conditional losses or loss bounds, but cannot produce an unbiased probability-weighted ECL by silently using equal weights. Return measurement\_unresolved for the ECL claim. Trade receivables and certain other assets use the separately applicable simplified approach; the general loan staging rule must not overwrite AASB 9.5.5.15--16.

\hypertarget{full-loan-and-allowance-roll-forward}{%
\section{Full loan and allowance roll-forward}\label{method-bank-methods-full-loan-and-allowance-roll-forward}}

Assume a fixed-rate non-POCI amortised-cost loan held to collect its contractual cash flows, originated for 100, with one terminal payment 121 two years later, original EIR 10\%, no fees, collateral, modifications, FX or taxes. Since no payment falls due before maturity, an earlier default in this example is a separately established qualitative credit event, such as financial distress under the admitted default definition; it is not a missed instalment. Article~52 of ASBE~22 requires a definition consistent with credit-risk management and consideration of qualitative indicators \citep{mofasbe22}. At origination the assumed lifetime scenarios are: default in the first twelve months with probability 0.1, default during the following year with probability 0.1, and no default with probability 0.8. Either default gives terminal collection 60.5. The first twelve-month ECL includes lifetime loss from the first scenario only: 0.1\ensuremath{\times}(121\ensuremath{-}60.5)/1.1\ensuremath{^2}=5; initial lifetime ECL would be 10. Debit loan 100/ credit cash 100; debit impairment 5/credit allowance 5. Gross 100, allowance 5, net 95.

\enlargethispage{\baselineskip}
At year-one accrual the loan has survived the first window and is not yet credit-impaired: debit gross loan 10/credit interest 10. Survival alone would give remaining default probability 0.1/0.9=1/9. A subsequent, separately admitted credit forecast revises that conditional probability to 0.2; it is an authored forecast update, not a Bayes calculation from survival alone. Its lifetime weights are 0.8/0.2, with a separately established SICR finding; ECL=0.2\ensuremath{\times}60.5/1.1=11. Debit impairment 6/ credit allowance 6. Gross 110, allowance 11, net 99. The ratio change alone is not the defined SICR test. Only twelve months remain, so twelve-month and lifetime ECL both equal 11 under this revised forecast. This amount cannot distinguish the two classifications.

Immediately afterward, a distinct, ordered credit-impairment event changes expected maturity collections to 60.5. This is a forecast mean, not a guaranteed or maximum recovery. Required allowance is 110\ensuremath{-}60.5/1.1=55; debit impairment 44/credit allowance 44. Net 55. ASBE~22 Article~59(5) measures this non-POCI credit-impaired asset by gross balance less discounted expected receipts. Articles~38 and 39(2) then make amortised cost, after the allowance, the basis for subsequent-period interest \citep{mofasbe22}.

The bank is assumed still to have reasonable prospects of fuller recovery before maturity. The separate finding that some contractual cash will not reasonably be recovered first arises at settlement on this branch. If that finding arose earlier, ASBE~22 Article~43 would require the corresponding earlier write-off. Expected collections alone do not justify postponing an already-required write-off.

In the next full year interest revenue is 55\ensuremath{\times}10\%=5.5. The gross contractual accretion is 11. One consistent decomposed posting debits gross loan 11, credits interest revenue 5.5 and allowance 5.5: gross 121, allowance 60.5, net 60.5. The allowance increase of 5.5 is attributed to discount unwind in its roll-forward, with no additional impairment expense in this posting. This reconciles gross and net-interest bases without counting the same expected loss again. Collect 60.5: debit cash/credit gross 60.5. When no reasonable expectation of remaining recovery is established, debit allowance 60.5/credit gross 60.5. Both balances reach zero. A mere delinquency is not that write-off finding. Subsequent unexpected cash recovery is recorded with its own recovery-income method and evidence; it does not restore an old live receivable by default.

\begin{longtable}[]{@{}lrrrrr@{}}
\toprule
\begin{minipage}[b]{0.11\columnwidth}\raggedright
Moment\strut
\end{minipage} & \begin{minipage}[b]{0.14\columnwidth}\raggedleft
Cash, opening 200\strut
\end{minipage} & \begin{minipage}[b]{0.14\columnwidth}\raggedleft
Gross loan\strut
\end{minipage} & \begin{minipage}[b]{0.14\columnwidth}\raggedleft
Allowance\strut
\end{minipage} & \begin{minipage}[b]{0.14\columnwidth}\raggedleft
Deposit liability\strut
\end{minipage} & \begin{minipage}[b]{0.14\columnwidth}\raggedleft
Equity, opening 20\strut
\end{minipage}\tabularnewline
\midrule
\endhead
\begin{minipage}[t]{0.11\columnwidth}\raggedright
Initial loan and allowance\strut
\end{minipage} & \begin{minipage}[t]{0.14\columnwidth}\raggedleft
100\strut
\end{minipage} & \begin{minipage}[t]{0.14\columnwidth}\raggedleft
100\strut
\end{minipage} & \begin{minipage}[t]{0.14\columnwidth}\raggedleft
5\strut
\end{minipage} & \begin{minipage}[t]{0.14\columnwidth}\raggedleft
180\strut
\end{minipage} & \begin{minipage}[t]{0.14\columnwidth}\raggedleft
15\strut
\end{minipage}\tabularnewline
\begin{minipage}[t]{0.11\columnwidth}\raggedright
Year-one accrual and SICR\strut
\end{minipage} & \begin{minipage}[t]{0.14\columnwidth}\raggedleft
100\strut
\end{minipage} & \begin{minipage}[t]{0.14\columnwidth}\raggedleft
110\strut
\end{minipage} & \begin{minipage}[t]{0.14\columnwidth}\raggedleft
11\strut
\end{minipage} & \begin{minipage}[t]{0.14\columnwidth}\raggedleft
180\strut
\end{minipage} & \begin{minipage}[t]{0.14\columnwidth}\raggedleft
19\strut
\end{minipage}\tabularnewline
\begin{minipage}[t]{0.11\columnwidth}\raggedright
Credit impairment\strut
\end{minipage} & \begin{minipage}[t]{0.14\columnwidth}\raggedleft
100\strut
\end{minipage} & \begin{minipage}[t]{0.14\columnwidth}\raggedleft
110\strut
\end{minipage} & \begin{minipage}[t]{0.14\columnwidth}\raggedleft
55\strut
\end{minipage} & \begin{minipage}[t]{0.14\columnwidth}\raggedleft
180\strut
\end{minipage} & \begin{minipage}[t]{0.14\columnwidth}\raggedleft
\ensuremath{-}25\strut
\end{minipage}\tabularnewline
\begin{minipage}[t]{0.11\columnwidth}\raggedright
Maturity accrual\strut
\end{minipage} & \begin{minipage}[t]{0.14\columnwidth}\raggedleft
100\strut
\end{minipage} & \begin{minipage}[t]{0.14\columnwidth}\raggedleft
121\strut
\end{minipage} & \begin{minipage}[t]{0.14\columnwidth}\raggedleft
60.5\strut
\end{minipage} & \begin{minipage}[t]{0.14\columnwidth}\raggedleft
180\strut
\end{minipage} & \begin{minipage}[t]{0.14\columnwidth}\raggedleft
\ensuremath{-}19.5\strut
\end{minipage}\tabularnewline
\begin{minipage}[t]{0.11\columnwidth}\raggedright
Collection and write-off\strut
\end{minipage} & \begin{minipage}[t]{0.14\columnwidth}\raggedleft
160.5\strut
\end{minipage} & \begin{minipage}[t]{0.14\columnwidth}\raggedleft
0\strut
\end{minipage} & \begin{minipage}[t]{0.14\columnwidth}\raggedleft
0\strut
\end{minipage} & \begin{minipage}[t]{0.14\columnwidth}\raggedleft
180\strut
\end{minipage} & \begin{minipage}[t]{0.14\columnwidth}\raggedleft
\ensuremath{-}19.5\strut
\end{minipage}\tabularnewline
\bottomrule
\end{longtable}

Total interest 15.5 less impairment 55 is loss 39.5, equal to 100 advanced minus 60.5 received. The negative equity state survives; no missing branch or capital plug restores solvency. These five dates jointly test allowance, recognition and bank balance-sheet arithmetic. The same loss cannot also be deducted directly from gross before the stated write-off.

For a non-derecognized loan modification, AASB 9.5.4.3 requires gross PV at original EIR and a modification result, then allowance and SICR reassessment under 5.5.12. For example gross 110 changed to a single future 110 at original 10\% gives revised gross 100 and modification loss 10 before any allowance change. ASBE~22 Articles~42 and 56 similarly separate remeasurement of a non-derecognised asset from its credit-risk reassessment \citep{mofasbe22}. This is separate from the financing method's issuer-liability ten-percent test. A derecognition decision requires its own asset conditions and cannot reuse that threshold.

\newpage
\hypertarget{funding-capital-and-decision-policies}{%
\section{Funding, capital and decision policies}\label{method-bank-methods-funding-capital-and-decision-policies}}

Liquidity simulation settles the dated net payment queue by venue/currency, respecting encumbrance, available cash, committed line terms and admissible actions. Haircut/eligibility and transferability come from a versioned policy; market value alone is not available liquidity. With cash 20, due withdrawal 30 and an admitted usable facility 8, even a full draw leaves shortfall 2. Record the default/unresolved settlement according to the event's legal model. Raising exactly 2 of new equity without an investor event is an invented financing source. Interbank contagion requires explicit counterparties and settlement priorities; it cannot be inferred from aggregate deposits.

A RegulatoryAssessment reads an accounting view and an independent jurisdiction, date, perimeter, eligibility, deductions and exposure-weight policy. It outputs capital components, RWA/exposure measures, liquidity measures and reconciliations with their distinct denominators. Illustratively, equity 20 minus an assumed regulatory deduction 2 yields capital 18; assumed RWA 125 gives 14.4\%. Those inputs are synthetic, not a Basel rule, national capital approval or default threshold. No binding capital-policy numbers have been inspected for a target bank here. Use the Basel/national source and actual bank facts before real admission.

Behavior under capital/liquidity constraints is a separate economic model. Specify feasible actions, objectives, information set and ties before solving; failure to find an action returns policy\_failure or default under the stated contract. Halaj/Behn sources can motivate candidate behavior but cannot calibrate it for an individual bank. A forecast experiment must compare simple fixed-book, contractual runoff and explicitly estimated policy baselines with time-safe data, before claiming a predictive improvement. CECL and other bases retain separate authority/method branches in the dated-authority discussion.

